% Part: first-order-logic
% Chapter: semantics
% Section: intro-semantics

\documentclass[../../../include/open-logic-section]{subfiles}

\begin{document}

\olfileid[hi]{fol}{syn}{its}

\olsection{परिचय}

अभिव्यक्तियों को अर्थ देना अर्थविज्ञान का विषय है। अर्थविज्ञान की केंद्रीय
धारणा किसी !!a{structure} में संतुष्टि की है। !!^a{structure} भाषा के
निर्माण-अवयवों को अर्थ देती है: !!a{domain} वस्तुओं का एक अरिक्त समुच्चय
है। परिमाणकों का निर्वचन इस प्रांत पर विचरण करने के रूप में किया जाता है,
!!{constant}s को प्रांत के अवयव दिए जाते हैं, !!{function}s को !!{domain}
से उसी में जाने वाले फलन दिए जाते हैं, और !!{predicate}s को !!{domain} पर
संबंध दिए जाते हैं। मूल शब्दावली के इन नियतनों के साथ !!{domain} मिलकर
!!a{structure} बनाता है। !!^{variable}, !!{formula}s में आ सकते हैं; अतः
अर्थविज्ञान देने के लिए हमें उन्हें भी !!{domain} के !!{element}s नियत करने
होते हैं---इसे चर-नियतन कहते हैं। अंततः संतुष्टि-संबंध इन सबको एक साथ जोड़ता
है। कोई !!^a{formula}, !!a{structure}~$\Struct{M}$ में
!!a{variable} नियतन~$s$ के सापेक्ष संतुष्ट हो सकता है; इसे
$\Sat{M}{!A}[s]$ लिखते हैं। इस संबंध की परिभाषा भी~$!A$ की संरचना पर आगमन
से दी जाती है। उदाहरणतः तार्किक संयोजकों की सत्यता-सारणियों का उपयोग करके
$(!A \land !B)$ की संतुष्टि को $!A$ और ~$!B$ की संतुष्टि (या असंतुष्टि) के
आधार पर परिभाषित किया जाता है। फिर पता चलता है कि यदि !!{formula}~$!A$ कोई
!!a{sentence} है, अर्थात् उसमें कोई मुक्त चर नहीं है, तो !!{variable}
नियतन अप्रासंगिक है। इसलिए हम सीधे कह सकते हैं कि कोई !!{sentence}, किसी
!!{structure} में संतुष्ट है (या नहीं)।

!!{sentence}s के लिए संतुष्टि-संबंध $\Sat{M}{!A}$ के आधार पर हम वैधता,
अर्थगत तात्पर्य और संतुष्टनीयता की मूल अर्थगत धारणाएँ परिभाषित कर सकते हैं।
कोई !!^a{sentence} वैध है, $\Entails !A$, यदि प्रत्येक !!{structure} उसे
संतुष्ट करती है। !!{sentence}s का कोई समुच्चय उससे तात्पर्य रखता है,
$\Gamma \Entails !A$, यदि~$\Gamma$ के सभी !!{sentence}s को संतुष्ट करने
वाली प्रत्येक !!{structure},~$!A$ को भी संतुष्ट करती है। और !!{sentence}s
का कोई समुच्चय संतुष्टनीय है यदि कोई !!{structure} उसके सभी !!{sentence}s
को एक साथ संतुष्ट करती है। चूँकि !!{formula} आगमनात्मक रूप से परिभाषित हैं,
और संतुष्टि भी !!{formula}s की संरचना पर आगमन से परिभाषित होती है, इसलिए हम
अपने अर्थविज्ञान के गुण सिद्ध करने तथा परिभाषित अर्थगत धारणाओं को परस्पर
संबंधित करने के लिए आगमन का उपयोग कर सकते हैं।

\end{document}
